\subsection{应用}\label{sec:2.5}
\subsubsection{Wasserstein 距离的估计}\label{sec:2.5.1}
目前为止，我们关注使用经验测度 $\mu_n$ 在 Wasserstein 距离下估计测度 $\mu$. 正如本章标题所示，我们常常也关心 Wasserstein 距离本身的估计. 事实上，这种距离在许多统计任务中居于核心地位. 例如，$1$-Wasserstein 距离以“推土机\index{earth mover’s distance@推土机距离（earth mover’s distance）|hyperpage}距离”（earth mover's distance）之名在机器学习中的早期应用之一，是在信息检索中衡量图像之间的距离\cite{ref276}. 其他直接的例子还有最近邻\cite{ref29,ref261}与回归\cite{ref72,ref145}等.

估计的目标是利用独立同分布数据 $X_1,\ldots,X_m\sim\mu$ 与 $Y_1,\ldots,Y_n\sim\nu$，构造 $W_1(\mu,\nu)$ 的估计量 $\widehat W$. 一个自然候选是代入\index{plug-in estimator@代入估计量（plug-in estimator）|hyperpage}估计量（plug-in estimator）$\widehat W:=W_1(\mu_m,\nu_n)$，其中 $\mu_m,\nu_n$ 为上述样本对应的经验测度. 利用三角不等式与\autoref{prop:2.1}，容易得到该估计量的性能界：当 $d\ge3$ 时，
$$\mathbb E|W_1(\mu_m,\nu_n)-W_1(\mu,\nu)|
\le\mathbb EW_1(\mu_m,\mu)+\mathbb EW_1(\nu_n,\nu)\lesssim(m\wedge n)^{-1/d}.$$
这个粗略的界一般而言是尖锐的. 此外，使用本章末尾给出的更一般结果 \eqref{eq:2.21}，当 $d>2p$ 时还可得到
$$\mathbb E|W_p(\mu_m,\nu_n)-W_p(\mu,\nu)|
\le\mathbb EW_p(\mu_m,\mu)+\mathbb EW_p(\nu_n,\nu)\lesssim(m\wedge n)^{-1/d}.$$
当 $p>1$ 时，只有 $\mu$ 与 $\nu$ 足够接近，这个界才是尖锐的. 事实上，若 $p>1$ 且 $W_p(\mu,\nu)>c>0$，就能得到更好的速率. 例如，$p=2$ 时，文献\cite{ref220}证明
$$\mathbb E|W_2(\mu_m,\nu_n)-W_2(\mu,\nu)|\lesssim(m\wedge n)^{-2/d}.$$
该文还证明，这个速率基本上是尖锐的. 这一改进虽显著，却也表明 Wasserstein 距离的估计仍受到维数灾难的影响.

\subsubsection{假设检验}\label{sec:2.5.2}
这些上界可以直接用于两个经典的非参数假设检验问题：拟合优度\index{goodness-of-fit testing@拟合优度检验（goodness-of-fit testing）|hyperpage}检验（goodness-of-fit testing）与双样本检验（two-sample testing，亦称同质性检验）.

先考虑拟合优度检验. 给定来自未知分布 $\mu$ 的独立同分布观测 $X_1,\ldots,X_n$，以及固定分布 $\mu^0$，目标是检验
$$H_0:\mu=\mu^0\qquad\text{对}\qquad H_1:\mu\ne\mu^0.$$
例如，可取 $\mu^0$ 为 $\mathbb R^d$ 上的标准高斯分布，或 $[0,1]^d$ 上的均匀分布. 当 $d=1$ 时，有许多拟合优度检验，例如针对连续分布的 Kolmogorov--Smirnov 检验. 对离散分布，$\chi^2$ 检验也是一种常见选择；参见\cite[第14章]{ref205}.

注意，这两个假设等价地可以写为
$$H_0:W_1(\mu,\mu^0)=0\qquad\text{对}\qquad H_1:W_1(\mu,\mu^0)>0.$$
为了研究这种设定的理论极限，可考虑该检验问题的一个定量版本：对某个 $\varepsilon>0$，假设为
$$H^0:W_1(\mu,\mu^0)=0\qquad\text{对}\qquad H_1:W_1(\mu,\mu^0)>\varepsilon.$$
于是问题是，分离量 $\varepsilon$ 必须多大，才能保证第一类与第二类错误的合计保持较小.

考虑一个简单检验：对某个满足下式的阈值 $T_n^\alpha$，只要 $W_1(\mu_n,\mu^0)>T_n^\alpha$，就在水平 $\alpha\in(0,1)$ 下拒绝原假设：
$$\mu^0[W_1(\mu_n,\mu^0)>T_n^\alpha]=\alpha.$$
若观测到 $X_i=x_i$，$i=1,\ldots,n$，则这种检验的 $p$ 值可以计算为
$$\mu^0[W_1(\mu_n,\mu^0)>W_1(\mu_n^{\mathrm{obs}},\mu^0)],$$
其中
$$\mu_n^{\mathrm{obs}}=\frac1n\sum_{i=1}^n\delta_{x_i}.$$
计算 $T_n^\alpha$ 和 $p$ 值，都需要理解原假设下 $W_1(\mu_n,\mu^0)$ 的实际分布. 上述结果实际上离这种精度还很远，因为我们只知道 $\mu^0$ 支撑于单位立方体时 $\mathbb E[W_1(\mu_n,\mu^0)]$ 的上界. 不过，这些界已足以描绘出 Wasserstein 距离用于多元拟合优度检验时一个颇令人失望的图景. 事实上，结果提示，为控制该检验的第一类错误，必须取 $T_n^\alpha\gtrsim n^{-1/d}$，因而要有合理的功效，需要分离量 $\varepsilon\gg n^{-1/d}$. 即使维数并不很高，这样的分离水平也相当大.

事实上，上述检验对这个问题并不最优；利用另一种方法，可以在更小的分离量 $n^{-2/d}$ 下实现一致检验. 但一个简单事实仍然成立：$W_1$ 的缓慢收敛阻碍了基于 Wasserstein 距离的拟合优度检验的使用.

对此现象的一种解释是，基于 $W_1$ 的检验试图对过多备择假设保持功效. 为便于讨论，设 $\mu^0$ 是 $\mathbb R^d$ 上的标准高斯分布. $1$-Wasserstein 距离不会区别那些不等于 $\mu^0$ 的分布：高斯分布、具有光滑密度的分布、密度不连续的分布，甚至离散分布. 检验 $\{W_1(\mu_n,\mu^0)>T_n^\alpha\}$ 试图检测所有这些分布，因而力量分散，对所有备择假设的功效都较低. 可以把这种行为与简单的参数检验作比较，例如 Wald 检验 $\{|\bar X_n|>\tau_n\}$，其中 $\bar X_n=\frac1n\sum_{i=1}^nX_i$；它只试图检测 $\mu$ 的均值是否不同于 $\mu^0$ 的均值. 显然，即使 $\mu$ 是与 $\mu^0$ 相当远的 Rademacher 分布，这个检验也不能检测出来；但它将全部力量集中在均值位移上，因此当这种位移出现时，就能非常精确地检测到.

维数灾难也出现在双样本检验中：观察来自 $\mu$ 的样本 $X_1,\ldots,X_m$ 与来自 $\nu$ 的样本 $Y_1,\ldots,Y_n$，目标是检验
$$H_0:\mu=\nu\qquad\text{对}\qquad H_1:W_1(\mu,\nu)>\varepsilon.$$
将相应经验分布记为
$$\mu_m=\frac1m\sum_{i=1}^m\delta_{X_i},\qquad\nu_n=\frac1n\sum_{j=1}^n\delta_{Y_j}.$$
在这个背景下，当 $W_1(\mu_m,\nu_n)$ 很大时拒绝原假设是自然的. 与拟合优度检验类似，这类检验若要达到合理功效，所需样本量会随维数呈指数增长.

本节的结论是，Wasserstein 距离受到维数灾难影响，因此不适合用于中等维数的统计应用. \autoref{sec:2.8}将介绍 Wasserstein 距离的多种正则化，它们避开维数灾难，并已成功用于大规模统计应用.
